\section{Conclusion}

This project examined online news popularity through a task-oriented data analytics framework. Rather than presenting the work as a collection of methods, the report organized the analysis around six specific tasks: popularity classification, share prediction, article grouping, engagement pattern discovery, formatting analysis, and publication timing investigation. This structure made the project objectives clearer and connected the technical analysis to practical questions about article engagement.

The strongest results came from the supervised tasks, where ensemble models such as Gradient Boosting and Random Forest consistently outperformed simpler alternatives. Feature engineering showed that keyword statistics, topic probabilities, self-reference measures, and sentiment-related variables are among the most important predictors of online popularity. At the same time, the study also showed that exact engagement prediction remains difficult and that not all analytical tasks yield equally strong outcomes.

The exploratory tasks added useful context by showing that the dataset has limited natural cluster separation and by identifying a small number of interpretable engagement patterns. The formatting analysis provided practical insight into how structural article features relate to engagement, while the publication timing task highlighted the importance of careful interpretation under class imbalance.

Overall, the report demonstrates that online news popularity can be studied effectively through a structured combination of predictive modeling, exploratory analysis, and feature-based interpretation. The Spark extension further shows how the same workflow can be adapted to scalable distributed environments \cite{zaharia2016spark}.
